\section{Abstract}
	\textbf{Machine Learning (ML)} techniques have become the predominant choice for predictive systems in software engineering. The problem arises as these models often lacks transparency. Making it difficult for developers and engineers to verify the result and ensure that the ML model reflects the real world.
	
	This project introduces a novel software design pattern, \textbf{Relational Event Prediction System (REPS)}, which seeks to address interpretability concerns by modelling event relationships through a graph structured system.
	
	REPS defines a prediction system composed of independent nodes that retrieves values from parent nodes. These nodes are configured using JSON with a set of unique keywords. Each node represent either a \textbf{Sensor} or \textbf{Event}, an Event can contain a user defined expression using C\# compiled in run-time or choose an existing implementation of a ML model, such as Support Vector Machines and Random Forest. This means REPS is flexible in its use-case areas.
	
	The implementation demonstrates that REPS is a low-complexity system requiring only 30 configuration keywords, while remaining expressive enough to support a variety of use cases. Runtime performance measurements indicate minimal CPU usage due to parallel node execution, but memory consumption remains relatively high. A simulator was also developed for testing and verification purposes.
	
	This project concludes that while REPS is promising alternative to traditional ML-based systems, further work and testing is needed to determine to what extend REPS can compete.

	%The software industry and research has increasingly relied on Machine Learning (ML) for making predictions and determinations, where developers would select the ML model that fits best with their use-case requirements. But current day ML's predictions can be hard to understand in certain contexts and developers can't be certain that their ML model correctly represents their use-case. Through this project I will be designing and implementing a new software pattern which I call Relational Event Prediction System (REPS), a graph based prediction system which can with their relation to one another make predictions in a way a user can better understand it. There will be a focus on how to construct it and how it can be used by the user. This project is exploratory in nature to determine the viability of REPS. This project managed to determine that there is validity to REPS but failed to determine to what degree, it managed to create a low complexity pattern with low technology knowledge requirement and only having 30 unique key-words. Measurements of REPS was performed and there is a high memory usage.